\chapter*{Введение} % * не простaвляет номер
\addcontentsline{toc}{chapter}{Введение} % вносим в содержaние


\textbf{Актуальность}. Регрессионное тестирование является неотъемлемой частью процесса разработки 
программного обеспечения: после каждого изменения кода необходимо убедиться, что существующая 
функциональность не нарушена. По мере роста проекта тестовый набор увеличивается экспоненциально, 
что делает полный повторный прогон тестов дорогостоящим и времязатратным. Приоритизация тестовых случаев 
(Test Case Prioritization, TCP) позволяет упорядочить выполнение тестов таким образом, чтобы дефекты 
обнаруживались как можно раньше, не исключая ни одного теста из набора. Применение эволюционных методов 
оптимизации, в частности генетических алгоритмов, открывает возможность учитывать одновременно несколько 
критериев качества при формировании порядка выполнения тестов.

\textbf{Предмет исследования}. Методы приоритизации тестовых случаев на основе генетического алгоритма с 
многокомпонентной фитнесс-функцией, включающей покрытие кода, оценку вероятности дефектов (fault-proneness) 
и структурное разнообразие тестов (diversification) на основе расстояния Жаккара.

\textbf{Цель исследования}. Разработка и экспериментальная оценка алгоритма приоритизации тестовых случаев, 
превосходящего базовый генетический алгоритм и жадный алгоритм дополнительного покрытия по метрикам APFD и 
APFDc на реальных наборах данных.

\textbf{Зaдaчи:}
\begin{enumerate}[1.]
	\item Изучить существующие методы приоритизации тестовых случаев и метрики оценки их эффективности.
	\item Разработать фитнесс-функцию генетического алгоритма, учитывающую покрытие кода, fault-proneness и структурное разнообразие тестов.
	\item Реализовать алгоритм на языке Java с поддержкой загрузки данных в форматах CSV и JSON.
	\item Провести сравнительный эксперимент на разных наборах данных.
	\item Проанализировать результаты и определить условия применимости предложенного подхода.
\end{enumerate} 
